% 第六章：内核实测

\section{内核实测} \label{sec:test}

本章只报能在真实内核上测到的东西：调度类回调是否被调用、决策是否落地、代价多大、代价来自哪里。所有数字可由 \texttt{scripts/run\_kernel\_repeats.sh} 重跑，原始 CSV 与串口日志归档在 \texttt{docs/validation/2026-11-13/kernel/repeat/}。

\subsection{测量设置}

测量设置见\autoref{tab:setup}。

\begin{table}[H]
\centering
\caption{实测环境与负载}
\label{tab:setup}
\small
\begin{tabular}{@{}p{3.0cm}p{10.0cm}@{}}
\toprule
项 & 取值 \\
\midrule
内核 & Linux 6.8，\texttt{arch/x86/boot/bzImage}，镜像 \#19（SHA256 \texttt{077638d4\ldots}） \\
机器 & QEMU 4 vCPU，\texttt{-accel tcg}，1\,GB 内存，保留 ACPI \\
负载 & 8 个任务 $\times$ 2\,000\,000 次迭代 = $1.6\times10^7$ 个工作单元，两策略完成量必须相同 \\
唤醒 & \texttt{--sleep-us 200}：把每个任务切成 64 段，制造 512 次阻塞—唤醒 \\
亲和性 & 主口径 \texttt{none}（可跨核，\texttt{select\_task\_rq} 才会被调用） \\
重复 & 每配置 5 次独立测量，每次含两次独立启动（先 CFS 后 Yat） \\
统计量 & 8 个任务墙钟之和的比值 $\rho=\sum W_{\text{yat}}/\sum W_{\text{cfs}}$，报中位数与全距 \\
\bottomrule
\end{tabular}
\end{table}

为什么不用均值而用中位数：TCG 是纯软件模拟，同一配置连跑三次的墙钟可差 28\%（\S\ref{sec:difficulties}），均值会被单次抖动主导；5 次重复的中位数与全距已经足够回答"是否成倍变慢"这个量级问题，不足以给出 $p$ 值，本文也不给。

三组配置各自回答一个问题：

\begin{enumerate}
  \item \textbf{none + 200\,$\mu$s}（主口径）：可跨核、有唤醒，AJLR 完整走一遍——问"代价是多少"。
  \item \textbf{shared + 200\,$\mu$s}：全部任务 pin 到单核，此时 \texttt{nr\_cpus\_allowed==1}，\texttt{core.c} 的包装层直接返回该核，\textbf{调度类回调一次都不会被调用}——用来把"选核的代价"从"接入调度类的代价"里剥出来。
  \item \textbf{none + 无唤醒}：只有 \texttt{WF\_FORK}，唤醒路径的决策不发生——用来验证代价是否随唤醒次数增长。
\end{enumerate}

\subsection{调度类是否被调用}

先回答一个更基本的问题：这一切真的跑起来了吗。\autoref{tab:counters} 是主口径下的计数器。

\begin{table}[H]
\centering
\caption{主口径计数器（镜像 \#19，none + 200\,$\mu$s）}
\label{tab:counters}
\small
\begin{tabular}{@{}p{4.4cm}rp{7.6cm}@{}}
\toprule
计数器 & 值 & 含义 \\
\midrule
\texttt{ttwu\_ok} & 512 & 成功唤醒次数（= 8 任务 $\times$ 64 段） \\
\texttt{select\_calls} & 487 & 进入 \texttt{select\_task\_rq\_yat\_casched()} 的次数（95.1\%） \\
\texttt{select\_sticky\_idle} & 216 & 目标核空闲，直接复用上次核（占 44.4\%） \\
\texttt{select\_decided} & 271 & 走完整 AJLR 决策 \\
\texttt{jobs\_placed} & 271 & 决策后真正落地的作业数 \\
\texttt{place\_hit} / \texttt{miss} / \texttt{pending} & 996 / 0 / 21 & 接管率 $996/1017=97.9\%$ \\
\texttt{enq\_calls} & 495 & 入队次数 \\
\texttt{enq\_cpus\_allowed} & 1980 & 累计可运行核数，平均 4.0 \\
\texttt{balance\_runs} & 757 & balance 钩子执行次数 \\
\texttt{pick\_calls} / \texttt{pick\_null\_desync} & 1024 / 0 & 选下一个任务 1024 次；desync 为 0 表示 \texttt{nr\_running} 维护正确 \\
\texttt{deq\_sleep} & 495 & 因睡眠出队，与 \texttt{ttwu} 对称 \\
\bottomrule
\end{tabular}
\end{table}

三个读数：

\textbf{回调确实被调用。}512 次唤醒里有 487 次进入调度类（其余 45 次走了 \texttt{ttwu} 的 wakelist 快路径，属于 \texttt{core.c} 的选择，不是本调度类能控制的）。这排除了"亲和性把回调锁死"这个最常见的假阳性——同一份脚本在 \texttt{shared} 下 \texttt{select\_calls} 为 0，正好是反面对照。

\textbf{决策确实落地。}271 次完整决策产生 271 次放置，\texttt{place\_hit}=996、\texttt{place\_miss}=0、\texttt{pending}=21，接管率 97.9\%。"算出来了但没落地"是这类实现最容易出的错，这里可以排除。

\texttt{sticky\_idle} 占 44.4\% 说明近一半唤醒根本不需要重算——目标核还空着，复用上次的核既省时间又保局部性。这个比例也是后续优化的入口（\S\ref{sec:conclusion}）。

\subsection{代价与归因}

\begin{table}[H]
\centering
\caption{三组配置的墙钟比值 $\rho$（Yat / CFS，5 次重复）}
\label{tab:cost}
\small
\begin{tabular}{@{}lcccc@{}}
\toprule
配置 & 中位数 & 最小 & 最大 & 读法 \\
\midrule
none + 200\,$\mu$s（主口径） & 2.00 & 1.61 & 2.51 & 完整链路 \\
shared + 200\,$\mu$s（回调不被调用） & 1.70 & 1.55 & 2.40 & 只有接入代价 \\
none + 无唤醒 & 1.10 & 0.73 & 1.36 & 无唤醒路径 \\
\bottomrule
\end{tabular}
\end{table}

\autoref{tab:cost} 的三行必须一起读，单看第一行会得出错误结论：

\begin{enumerate}
  \item \textbf{主口径下 Yat 比 CFS 慢约一倍}（中位数 2.00）。这是本文最重要也最不好看的数字，直接说明当前实现还不具备可用性。
  \item \textbf{把回调关掉，仍然慢}（1.70）。\texttt{shared} 下 \texttt{select\_calls}=0，AJLR 一次都没跑，慢的原因只能来自"任务进了这个调度类"本身——入队出队、\texttt{nr\_running} 维护、balance 钩子。这一条排除了"是 AJLR 选核太贵"的解释。
  \item \textbf{没有唤醒时，两者没有可分辨的差异}（1.10，全距 $[0.73,1.36]$ 跨 1）。此时只剩 \texttt{WF\_FORK}，唤醒相关的回调不发生，基础调度机制（时间片、pick、上下文切换）不构成代价。这一行同时排除一种解释：差距不是“Yat 这个调度类”的固定开销，而是随唤醒次数增长的。
\end{enumerate}

把三条合起来，代价的归属是清楚的：\textbf{主体不在选核算法，而在“唤醒时进入这个调度类”这条路径上}。1.70 与 2.00 之差是选核的边际代价（约 $2.00/1.70-1\approx18\%$）；1.10 与 1.70 之差是唤醒路径的代价，占主要部分。计数器也支持这个判断——\texttt{balance\_runs}=757 而唤醒只有 512 次，意味着几乎每个调度点都在执行 balance，而 balance 的实现在最坏情况下要扫描全部 DAG 的就绪池并逐核比较。

需要克制的一点：镜像 \#18（修复前）单次测得 2.295，\#19 的 5 次中位数是 2.00（区间 $[1.61,2.51]$），该差值完全落在噪声内，因此\textbf{不能}说"锁修复提升了性能"。能说的是"锁修复消除了 \texttt{lockdep} 告警，且没有带来性能回归"。

\subsection{前期工作：仿真器}

项目前期实现过一个 Java 17 模拟器，用同一套缓存模型比较 WFD 与缓存感知启发式，并用 5 个种子 $\times$ 100 个用例给出统计区间。它得到的结论是：在理想化模型下，旧启发式相对 WFD 没有稳定收益（配对均值 $-3.52\%$，区间整段在零以下），把"任务上次在哪个核跑过"接回评分函数后能到 $+0.99\%$ 与 $+6.79$ 个百分点。

这些数字\textbf{不作为}本文的证据：模拟器的执行时间与命中判定都由自建模型给出，不含抢占、迁移、TLB 与真实 LLC 竞争，口径稍一变动结论就翻转（同一批数据在三套口径下分别是 11.66\%、$-0.95\%$、$-3.52\%$，详见 \texttt{docs/validation/}）。它留下的可用资产是"设计期该验证什么"的清单，代码与全部实验结果仍在仓库中，供需要复现的人查阅。

\subsection{有效性与边界}

\begin{itemize}
  \item \textbf{平台是 QEMU/TCG，不是裸金属}。TCG 无硬件虚拟化，墙钟与真实机器不可比；本文的结论只到"在 TCG 下成倍变慢"这个量级，迁移到裸机需要重测。
  \item \textbf{单负载形状}：8 任务、200\,$\mu$s 睡眠、64 段。结论对"频繁唤醒 + 可跨核"的负载成立，对纯计算或强亲和性负载不适用（后者回调根本不被调用）。
  \item \textbf{5 次重复只够给区间}。同一配置的三次运行可差 28\%，本文因此报中位数与全距，不报 $p$ 值，也不把 5\% 量级的差异当结论。
  \item \textbf{没有真实 LLC 竞争}。QEMU 的 4 vCPU 共享宿主机缓存，guest 内的缓存层级是模拟的；本文不声称测到了缓存收益，只测到调度开销。
  \item \textbf{亲和性语义}：\texttt{shared} 配置下 \texttt{select\_task\_rq} 不被调用，这类数据只能谈时间片分配，不能谈 AJLR——这一点在脚本头注释与第\ref{sec:difficulties}节都写明了，避免以后误读。
\end{itemize}